%%%%%%%%%%%%%%%%%%%%%%%%%%%%%%%%%%%%%%%%%%%%%%%%%%%%
\section{関連研究}

\subsection{頻出アイテムセットマイニング}

Apriori~\cite{agrawal1994fast} と FP-Growth~\cite{han2000mining} は頻出アイテムセットマイニング（FIM）の
基盤となる手法である。
ECLAT~\cite{zaki1997new} は垂直データ形式を用いた代替手法を提供する。
FIM の包括的な調査は~\cite{fournier2017survey,luna2019fim,fournierviger2022survey} に詳しい。
周期的頻出パターンマイニング PFPM~\cite{tanbeer2009pfpm} や，
複数の最小サポートを導入して希少アイテム問題~\cite{liu1999multiple} に対処する研究も存在するが，
いずれもグローバルな評価に依存する点は変わらない。
これらの手法はいずれもデータセット全体でサポートを評価するため、特定の時間区間でのみ高頻度で出現する
アイテムセットを検出できない。

\subsection{密集区間マイニング}

パターンとその密集区間を同時に出力する手法は2種類に大別される。
\textit{発生ギャップ型}手法には、LPFIM~\cite{yin2014lpfim}、LPPM~\cite{saleh2011lppm}、
RPM~\cite{kiran2017rpm} がある。
LPFIM は出現ギャップが閾値に達したとき区間を分割し、LPPM はノイズロバストな検出のために
累積スピルオーバーを導入し、RPM は最大周期内のギャップで連続する出現をグループ化する。
これらの手法はギャップ制約や周期閾値を通じて密度を間接的に推定するため、パラメータ調整が困難であり、
直接的なスライディングウィンドウ頻度に対して結果が近似的となる。
\textit{直接型}手法としては、固定時間粒度でパターンをマイニングするため任意幅の区間を特定できない
GLFMiner~\cite{fournierviger2021glf} と、アイテムセットごとに最適な単一区間のみを返す
SIM~\cite{kiran2015sim} がある。
これらの手法はいずれもストリーミングをサポートせず、スキャン頻度を制御する機構を持たない。

\subsection{ストリームマイニング}

スライディングウィンドウを用いたストリームマイニングはインクリメンタルFIM として研究されており
~\cite{lee2001sliding,leung2006dstree}、ウィンドウがサポート計算の時間的スコープを制限する。
しかしこれらの手法はウィンドウ内のグローバルサポートを報告するものであり、
すべての密集部分区間を列挙しない。
バースト検出~\cite{kleinberg2003bursty,yang2021burstsketch} はストリーム中の
頻度スパイクを検出するが、アイテムセット共起の密集区間は列挙しない。

\subsection{本研究の立ち位置}

密集パターンの本質は、あるアイテムセットが特定の連続区間内でのみ局所的に高頻度で出現する点にある。
本研究では、スライディングウィンドウを用いてウィンドウ内の出現頻度を直接評価することで、
発生ギャップ制約に依存せずに局所密度を定量化する。
これによりパラメータ感度問題を原理的に回避できる。

既存の手法群を整理すると、(1) パターンのみを出力するFIM系手法、(2) パターンと区間を同時に出力するが
ギャップ制約経由で密度を間接推定するLPFIM/LPPM/RPM、(3) 直接的な局所頻度を評価するが
区間の粒度や個数に制限があるGLFMiner/SIM、(4) ストリームマイニングではあるが密集部分区間を
列挙しないインクリメンタルFIM、のいずれも「スライディングウィンドウによる直接的な局所頻度評価」と
「アイテムセットの反単調性を組み合わせた密集区間の網羅的列挙」を
ストリーミング環境で実現するものではない。
本稿はこの方向性をストリーミング対応アルゴリズムとして定式化・実装する試みである。
